\documentclass[11pt,letterpaper]{article}

% Standalone source: upload this file to Overleaf and compile with pdfLaTeX.
\usepackage[T1]{fontenc}
\usepackage[utf8]{inputenc}
\usepackage{charter}
\usepackage[margin=1in]{geometry}
\usepackage{microtype}
\usepackage[hidelinks]{hyperref}
\hypersetup{
  pdftitle={Research Statement},
  pdfauthor={Mst Afsana Mimi},
  pdfsubject={PhD application, University of Arkansas, Fall 2027}
}
\input{glyphtounicode}
\pdfgentounicode=1
\pagestyle{empty}
\setlength{\parindent}{0pt}
\setlength{\parskip}{0.75\baselineskip}
\linespread{1.04}
\frenchspacing

\begin{document}

\begin{center}
  {\Large\bfseries Research Statement}\\[6pt]
  {\large Mst Afsana Mimi}\\[4pt]
  {\small Prospective PhD Student, Fall 2027\enspace$\vert$\enspace University of Arkansas}
\end{center}
\vspace{2pt}

My research interests are in statistical machine learning, particularly how AI systems choose and use information to make predictions. I came to these questions through my training in statistics and my research on high-dimensional data, fairness in prediction, and, more recently, agentic AI. In my current work with relational databases, I am interested in how a system can decide which information is useful for a prediction and how to bring that information together.

I recently completed my master's degree in Statistics at the University of Massachusetts Amherst, where I now work as a Research Assistant in the DREAM Lab under the supervision of Prof.\ Marco Serafini. My current project focuses on agentic AI for relational prediction. The goal is to develop a compound AI system that can answer predictive queries over relational databases. We use database profiling statistics to guide how attributes are aggregated across meta-paths. This work has made me think more closely about the decisions that come before a prediction: which relationships to follow, which information to select, and how to use it.

During my master's studies, I worked on Gaussian process regression with fairness constraints. The goal was to improve fairness in predictions while maintaining predictive accuracy. Before that, at the University of Dhaka, I studied variable selection for high-dimensional censored data. I developed a feature-selection method and evaluated it through simulation studies and applications to biomedical data. This research was published in \emph{Statistics in Medicine} in 2021. These projects gave me experience with different modeling problems and helped shape my interest in how the choice of methods and variables affects a prediction.

For my PhD, I would like to build on this experience by studying Bayesian optimization and search and planning methods for language models. My work with Gaussian process regression provides a starting point for learning how probabilistic models can guide exploration when observations are expensive or limited. I also want to study how agentic AI systems decide what information to retrieve and what action to take next. I see a close connection to my current research: in both settings, a system must choose which information to use before it can make a useful prediction or decision.

The Arkansas Machine Learning Group interests me because of its research on Bayesian optimization, active learning, and planning for language models. I would like to explore these topics in greater depth while bringing my background in statistical modeling and my current experience with relational prediction to the group's research. Through doctoral study, I hope to develop the methods and research skills needed to investigate how learning systems can make better decisions with the information available to them.

\end{document}
